% ==============================================================================
% CONFIGURAZIONE DEI PACCHETTI (PACKAGES) - FISICA GENERALE 1
% ==============================================================================

\usepackage[utf8]{inputenc}
\usepackage[T1]{fontenc}
\usepackage[italian]{babel}
\usepackage{lmodern}
\usepackage{microtype}

% --- Geometria e layout della pagina ---
\usepackage[
    a4paper,
    top=2.5cm,
    bottom=2.5cm,
    left=2.7cm,
    right=2.7cm,
    headheight=26pt,
    footskip=1.2cm
]{geometry}

% --- Pacchetti matematici e fisici avanzati ---
\usepackage{amsmath}
\usepackage{amssymb}
\usepackage{amsthm}
\usepackage{mathtools}
\usepackage{amsfonts}
\usepackage{mathrsfs}
\usepackage{bm}
\usepackage{cancel}
\usepackage{siunitx}
\usepackage{esint}

% Configurazione siunitx per italiano
\sisetup{
    output-decimal-marker = {,},
    separate-uncertainty = true,
    per-mode = symbol,
    range-phrase = { $\div$ },
    range-units = single
}
\DeclareSIUnit{\ly}{ly}
\DeclareSIUnit{\parsec}{pc}
\DeclareSIUnit{\astronomicalunit}{AU}
\DeclareSIUnit{\tonne}{t}
\DeclareSIUnit{\calorie}{cal}
\DeclareSIUnit{\atmosphere}{atm}
\DeclareSIUnit{\torr}{Torr}
\DeclareSIUnit{\bar}{bar}

% --- Gestione grafica, TikZ e grafici ---
\usepackage{graphicx}
\usepackage{xcolor}
\usepackage{tikz}
\usetikzlibrary{
    calc,
    arrows.meta,
    positioning,
    patterns,
    decorations.pathmorphing,
    decorations.markings,
    backgrounds,
    shapes.geometric,
    angles,
    quotes
}
\usepackage{pgfplots}
\pgfplotsset{compat=1.18}

% --- Strutture, tabelle e liste ---
\usepackage{booktabs}
\usepackage{tabularx}
\usepackage{array}
\usepackage{multirow}
\usepackage{enumitem}
\setlist{itemsep=0.2em, topsep=0.3em}

% --- Box decorati con tcolorbox ---
\usepackage[most]{tcolorbox}

% --- Intestazioni e piè di pagina ---
\usepackage{fancyhdr}
\pagestyle{fancy}
\fancyhf{}
\fancyhead[LE,RO]{\thepage}
\fancyhead[RE]{\nouppercase{\leftmark}}
\fancyhead[LO]{\nouppercase{\rightmark}}
\renewcommand{\headrulewidth}{0.4pt}
\renewcommand{\footrulewidth}{0pt}

% Stile per inizio capitoli
\fancypagestyle{plain}{
    \fancyhf{}
    \fancyfoot[C]{\thepage}
    \renewcommand{\headrulewidth}{0pt}
}

% --- Titoli dei capitoli e sezioni eleganti ---
\usepackage{titlesec}
\titleformat{\chapter}[display]
  {\normalfont\sffamily\huge\bfseries\color{navyblue}}
  {\flushright\normalsize\scshape\color{gray} CAPITOLO \thechapter}
  {0.5ex}
  {\titlerule[1.5pt]\vspace{1ex}}
  [\vspace{1ex}\titlerule]

\titleformat{\section}
  {\normalfont\sffamily\Large\bfseries\color{navyblue}}
  {\thesection}{1em}{}

\titleformat{\subsection}
  {\normalfont\sffamily\large\bfseries\color{forestgreen}}
  {\thesubsection}{1em}{}

\titleformat{\subsubsection}
  {\normalfont\sffamily\normalsize\bfseries\color{charcoalgray}}
  {\thesubsubsection}{1em}{}

% --- Collegamenti ipertestuali e riferimenti intelligenti ---
\usepackage{hyperref}
\hypersetup{
    colorlinks=true,
    linkcolor=navyblue,
    citecolor=forestgreen,
    urlcolor=purpleaccent,
    filecolor=navyblue,
    pdfauthor={Nicola Beccaceci},
    pdftitle={Fisica Generale 1: Appunti Completi, Teoria, Dimostrazioni ed Esercizi},
    pdfsubject={Fisica Generale 1 - Meccanica, Fluidi, Termodinamica},
    pdfkeywords={Fisica 1, Meccanica, Dinamica, Gravitazione, Fluidi, Termodinamica, Ingegneria Gestionale}
}
\usepackage{cleveref}

% --- Utility aggiuntive ---
\usepackage{caption}
\captionsetup{font=small, labelfont=bf, labelsep=period}
\usepackage{subcaption}
\usepackage{pdfpages}
